\documentclass{article}
\usepackage{eartharxiv}
\usepackage[numbers,square]{natbib}
\usepackage[utf8]{inputenc}
\usepackage[T1]{fontenc}
\usepackage{xcolor}
\usepackage{lineno}
\usepackage{graphicx}
\usepackage{booktabs}
\usepackage{hyperref}
\usepackage{authblk}
\hypersetup{colorlinks=true,linkcolor=purple,urlcolor=blue,citecolor=cyan,anchorcolor=black}

\title{PAG ETR · um modelo lógico e reprodutível de analogia geológica para priorização de investigação geoquímica de elementos terras raras em territórios com dados incompletos. Aplicação em Mato Grosso do Sul, Brasil}
\author[1]{Carlos Busón Buesa\thanks{\texttt{carlos.buson@ufms.br}}}
\affil[1]{Universidade Federal de Mato Grosso do Sul · Programa de Pós-Graduação em Tecnologias Ambientais · FAENG}
\date{}

\begin{document}
\maketitle

\begin{abstract}
\textit{Resumo a redigir após a auditoria editorial final. O pacote não preenche esta seção por inferência.}
\end{abstract}

\noindent\textbf{Palavras-chave} elementos terras raras, prospecção geoquímica, analogia geológica, reprodutibilidade, Mato Grosso do Sul

\section{Introdução}
Conteúdo em elaboração a partir do manuscrito mestre preservado no repositório.

\section{Contexto geológico e metalogenético}
Conteúdo em elaboração a partir do manuscrito mestre preservado no repositório.

\section{Dados e proveniência}
O pacote reprodutível preserva as fontes congeladas usadas no piloto espacial, as evidências derivadas, as três malhas de análise, os resultados congelados, os resultados reproduzidos e os registros de auditoria.

\section{Metodologia PAG ETR}
PAG ETR usa estados lógicos categóricos. No piloto espacial materializado, N0 representa condição não avaliável e N1 representa modelo aberto. A ausência de evidência materializada não é convertida em ausência geológica.

\subsection{Biblioteca de modelos M1 a M11}
A especificação executável e o catálogo de testes estão preservados no diretório \texttt{methodology}.

\subsection{Arquitetura multiescala}
As malhas de 250, 500 e 1000 km$^2$ são avaliadas de forma independente.

\section{Implementação do piloto espacial em Mato Grosso do Sul}
O piloto espacial reproduz M2 a partir do Corpo Feixe dos Morros e M4 a partir da Formação Bocaina, Formação Tamengo e registros de fósforo preservados nas bases congeladas do pacote.

\section{Resultados}
A reprodução validada recupera 49 células em estado de modelo aberto na malha de 250 km$^2$, 37 na malha de 500 km$^2$ e 24 na malha de 1000 km$^2$.

\begin{table}[ht]
\centering
\caption{Resultado congelado e reproduzido do piloto espacial M2 e M4}
\begin{tabular}{rrrrr}
\toprule
Escala km$^2$ & Células & N1 & M2 & M4 \\
\midrule
250 & 1554 & 49 & 3 & 46 \\
500 & 793 & 37 & 3 & 34 \\
1000 & 412 & 24 & 2 & 22 \\
\bottomrule
\end{tabular}
\end{table}

\begin{figure}[ht]
\centering
\includegraphics[width=0.92\linewidth]{figures/PAG_ETR_250km2.png}
\caption{PAG ETR na malha de 250 km$^2$. Figura preservada no pacote reprodutível.}
\end{figure}

\section{Validação e testes}
A validação compara os identificadores exatos dos hexágonos reproduzidos com os resultados congelados. O log V2 registra PASS nas três escalas.

\section{Discussão}
Conteúdo em elaboração a partir do manuscrito mestre preservado no repositório.

\section{Limitações}
O pacote reproduz o piloto espacial materializado M2 e M4. Ele não reconstrói por inferência os modelos que permaneceram sem requisitos ou geometrias espacializadas suficientes no corte congelado. O resultado PAG ETR identifica zonas de investigação e não demonstra jazimento, recurso, reserva, teor econômico ou autorização de campo.

\section{Conclusões}
Conteúdo em elaboração a partir do manuscrito mestre preservado no repositório.

\section*{Disponibilidade de código e dados}
O repositório inclui scripts de reprodução, testes, dados congelados, resultados reproduzidos, auditorias, figuras e a interface cartográfica publicada por GitHub Pages.

\end{document}
